% Research paper scaffold — source: ~/.agents/paper-template/ (writepaper_project trigger).
% Build: bash build.sh   (latexmk -> main.pdf; texlive-full is installed, no package is missing)
\documentclass[11pt,a4paper]{article}

\usepackage[utf8]{inputenc}
\usepackage[T1]{fontenc}
\usepackage{lmodern}
\usepackage[margin=1in]{geometry}
\usepackage{microtype}

% --- math ---------------------------------------------------------------
\usepackage{amsmath,amssymb,amsthm,mathtools}
\usepackage{bm}
\usepackage{siunitx}

% --- floats, tables, figures --------------------------------------------
\usepackage{graphicx}
\usepackage{xcolor}
\usepackage{booktabs}
\usepackage{tabularx}
\usepackage{subcaption}
\usepackage{tikz}
\usepackage{pgfplots}
\pgfplotsset{compat=1.18}

% --- algorithms ---------------------------------------------------------
\usepackage[ruled,vlined,linesnumbered]{algorithm2e}

% --- bibliography -------------------------------------------------------
\usepackage[backend=biber,style=numeric-comp,sorting=none,maxbibnames=99]{biblatex}
\addbibresource{refs.bib}

% --- links (load late) --------------------------------------------------
\usepackage[colorlinks=true,allcolors=blue]{hyperref}
\usepackage[nameinlink,capitalise]{cleveref}

% --- theorem environments ----------------------------------------------
\theoremstyle{plain}
\newtheorem{theorem}{Theorem}[section]
\newtheorem{lemma}[theorem]{Lemma}
\newtheorem{proposition}[theorem]{Proposition}
\newtheorem{corollary}[theorem]{Corollary}
\newtheorem{conjecture}[theorem]{Conjecture}
\theoremstyle{definition}
\newtheorem{definition}[theorem]{Definition}
\newtheorem{assumption}[theorem]{Assumption}
\newtheorem{example}[theorem]{Example}
\theoremstyle{remark}
\newtheorem{remark}[theorem]{Remark}

% Every gap is marked, never guessed. Grep \TODO before submitting.
\newcommand{\TODO}[1]{\textcolor{red}{\textbf{[TODO: #1]}}}

\title{\TODO{title}}
\author{Brusk Kawa Abdalla\thanks{\href{mailto:math@brwsk.xyz}{math@brwsk.xyz}}}
\date{\today}

\begin{document}
\maketitle

\begin{abstract}
\TODO{abstract: problem, method, main formal result, main empirical result, significance}
\end{abstract}

\noindent\textbf{Keywords:} \TODO{keywords}

\section{Introduction}
\label{sec:intro}
\TODO{motivation, the gap in existing work, and what this paper does}

\paragraph{Contributions.}
\begin{itemize}
  \item \TODO{contribution 1 — formal}
  \item \TODO{contribution 2 — algorithmic}
  \item \TODO{contribution 3 — empirical}
\end{itemize}

\section{Related work}
\label{sec:related}
\TODO{real, checkable references only — no invented citations}

\section{Preliminaries and notation}
\label{sec:prelim}
\TODO{define every symbol before use; keep a notation table}

\begin{table}[t]
  \centering
  \caption{Notation.}
  \label{tab:notation}
  \begin{tabular}{ll}
    \toprule
    Symbol & Meaning \\
    \midrule
    \TODO{sym} & \TODO{meaning} \\
    \bottomrule
  \end{tabular}
\end{table}

\section{Problem statement}
\label{sec:problem}
\begin{assumption}\label{as:main}
\TODO{state assumptions explicitly}
\end{assumption}

\section{Method}
\label{sec:method}
\TODO{the construction, derived rather than asserted}

\section{Theoretical results}
\label{sec:theory}

\begin{definition}\label{def:main}
\TODO{definition}
\end{definition}

\begin{lemma}\label{lem:main}
\TODO{statement}
\end{lemma}
\begin{proof}
\TODO{proof, or a sketch with the full argument in \Cref{app:proofs}}
\end{proof}

\begin{theorem}\label{thm:main}
\TODO{statement}
\end{theorem}
\begin{proof}
\TODO{proof}
\end{proof}

\subsection{Complexity}
\label{sec:complexity}
\TODO{time / space / sample complexity or numerical error bounds}

\section{Algorithm}
\label{sec:algorithm}
\begin{algorithm}[t]
  \caption{\TODO{name}}
  \label{alg:main}
  \KwIn{\TODO{input}}
  \KwOut{\TODO{output}}
  \TODO{steps}
\end{algorithm}

\section{Implementation}
\label{sec:implementation}
\TODO{architecture and the design decisions from docs/DECISIONS.md}

\section{Experimental setup}
\label{sec:setup}
\TODO{hardware, software versions, seeds, datasets, hyperparameters — enough to rerun}

\section{Results}
\label{sec:results}
\TODO{every number here comes from a real run; report n, mean $\pm$ CI, test statistic, p-value, effect size}

\begin{table}[t]
  \centering
  \caption{\TODO{caption}}
  \label{tab:results}
  \begin{tabular}{lccc}
    \toprule
    Method & $n$ & Mean $\pm$ 95\% CI & Effect size \\
    \midrule
    \TODO{row} & & & \\
    \bottomrule
  \end{tabular}
\end{table}

\subsection{Ablations}
\label{sec:ablations}
\TODO{what each component contributes}

\section{Discussion}
\label{sec:discussion}
\TODO{what the results mean, and what they do not show}

\section{Limitations and threats to validity}
\label{sec:limitations}
\TODO{scope conditions, failure cases, confounds}

\section{Conclusion and future work}
\label{sec:conclusion}
\TODO{conclusion}

\printbibliography

\appendix

\section{Full proofs}
\label{app:proofs}
\TODO{complete arguments deferred from \Cref{sec:theory}}

\section{Reproducibility}
\label{app:repro}
\TODO{exact commands, commit hash, environment}

\end{document}
